\documentclass[11pt,letterpaper,openany]{book}
\usepackage[utf8]{inputenc}
\usepackage[T1]{fontenc}
\usepackage[english]{babel}
\usepackage{amsmath}
\usepackage{amssymb}
\usepackage{amsthm}
\usepackage{graphicx}
\usepackage{float}
\usepackage{geometry}
\usepackage{fancyhdr}
\usepackage{listings}
\usepackage{xcolor}
\usepackage{tikz}
\usepackage{pgfplots}
\usepackage{hyperref}
\hypersetup{
    colorlinks=true,
    linkcolor=blue,
    filecolor=magenta,
    urlcolor=cyan,
    citecolor=red,
    pdftitle={QntLab2026 Cookbook: Complete Guide to Quantitative Trading},
    pdfauthor={Q23 Development Team},
    pdfsubject={Quantitative Trading, Systematic Strategies, Factor Investing},
    pdfkeywords={quantitative trading, systematic strategies, factor investing, risk management, portfolio optimization},
    pdfcreator={LaTeX with hyperref},
    bookmarks=true,
    bookmarksopen=true,
    bookmarksnumbered=true,
    breaklinks=true
}
\usepackage{booktabs}
\usepackage{multirow}
\usepackage{subcaption}
\usepackage{enumitem}
\usepackage{algorithm}
\usepackage{algpseudocode}
\usepackage{tcolorbox}
\usepackage{etoolbox}

% Page geometry
\geometry{margin=1in,top=1.25in,bottom=1.25in}
\linespread{1.1} % Slightly better line spacing

% Colors for code
\definecolor{codegreen}{rgb}{0,0.6,0}
\definecolor{codegray}{rgb}{0.5,0.5,0.5}
\definecolor{codepurple}{rgb}{0.58,0,0.82}
\definecolor{backcolour}{rgb}{0.95,0.95,0.92}

% Code listings setup
\lstdefinestyle{mystyle}{
    backgroundcolor=\color{backcolour},
    commentstyle=\color{codegreen},
    keywordstyle=\color{magenta},
    numberstyle=\tiny\color{codegray},
    stringstyle=\color{codepurple},
    basicstyle=\ttfamily\footnotesize,
    breakatwhitespace=false,
    breaklines=true,
    captionpos=b,
    keepspaces=true,
    numbers=left,
    numbersep=5pt,
    showspaces=false,
    showstringspaces=false,
    showtabs=false,
    tabsize=2
}
\lstset{style=mystyle}

% Headers and footers
\pagestyle{fancy}
\fancyhf{}
\fancyhead[L]{\leftmark}
\fancyhead[R]{\thepage}
\fancyfoot[C]{\textit{QntLab2026 Cookbook -- \today}}
\renewcommand{\headrulewidth}{0.4pt}
\renewcommand{\footrulewidth}{0.4pt}

% Custom theorem environments
\newtheorem{theorem}{Theorem}[chapter]
\newtheorem{lemma}[theorem]{Lemma}
\newtheorem{proposition}[theorem]{Proposition}
\newtheorem{corollary}[theorem]{Corollary}
\newtheorem{definition}[theorem]{Definition}
\newtheorem{example}[theorem]{Example}
\newtheorem{remark}[theorem]{Remark}

% Recipe environment
% \newtcolorbox[auto counter]{recipe}[1][]{
%     colback=blue!5!white,
%     colframe=blue!75!black,
%     title=Recipe~\thetcbcounter: #1,
%     fonttitle=\bfseries,
%     sharp corners,
%     boxrule=0.5mm,
% }

% Key insight box
\newtcolorbox{keyinsight}{
    colback=green!5!white,
    colframe=green!75!black,
    title=Key Insight,
    fonttitle=\bfseries,
    sharp corners,
    boxrule=0.5mm,
}

% Production note box
\newtcolorbox{productionnote}{
    colback=red!5!white,
    colframe=red!75!black,
    title=Production Note,
    fonttitle=\bfseries,
    sharp corners,
    boxrule=0.5mm,
}

% Custom commands
\newcommand{\R}{\mathbb{R}}
\newcommand{\N}{\mathbb{N}}
\newcommand{\Z}{\mathbb{Z}}
\newcommand{\Q}{\mathbb{Q}}
\newcommand{\E}{\mathbb{E}}
\newcommand{\Var}{\mathrm{Var}}
\newcommand{\Cov}{\mathrm{Cov}}
\newcommand{\Corr}{\mathrm{Corr}}
\newcommand{\Prob}{\mathbb{P}}
\newcommand{\IC}{\mathrm{IC}}
\newcommand{\IR}{\mathrm{IR}}
\newcommand{\Sharpe}{\mathrm{Sharpe}}
\newcommand{\MaxDD}{\mathrm{MaxDD}}
\newcommand{\OFI}{\mathrm{OFI}}
\newcommand{\VPIN}{\mathrm{VPIN}}
\newcommand{\GLFT}{\mathrm{GLFT}}

\title{QntLab2026 Cookbook: \\
The Complete Guide to Quantitative Trading \\
\textit{From Theory to Production Mastery}}


\author{Q23 Development Team \\
\textit{Inspired by Max Dama, Powered by Battle-Tested Systems}}

\date{2026}

\begin{document}

% Prevent empty pages in chapters
\makeatletter
\@openrightfalse
\makeatother

\maketitle

\setcounter{tocdepth}{3}
\tableofcontents

\chapter*{Foreword: Cooking Up Quant Alpha}
\addcontentsline{toc}{chapter}{Foreword}

Welcome to \textit{QntLab2026 Cookbook}, the most comprehensive guide to quantitative trading ever assembled. This is not just another academic text or superficial tutorial—it's a production-grade cookbook that takes you from basic concepts to building and deploying world-class systematic strategies.

\textbf{Why a Cookbook?}

In the kitchen of quantitative finance, recipes are everything. A great chef doesn't just understand ingredients—they know how to combine them, when to apply heat, and how to avoid burning the dish. Similarly, successful quant traders don't just understand factors—they know how to combine them, when to apply leverage, and how to avoid blowing up portfolios.

\textbf{What Makes This Different?}

\begin{itemize}
    \item \textbf{Production-Ready Code:} Every recipe comes with working Q23 system code
    \item \textbf{Real-World Examples:} Strategies that have been live in production
    \item \textbf{Risk Management First:} Every recipe emphasizes robustness and drawdown control
    \item \textbf{Complete Stack:} From factor discovery to portfolio deployment
    \item \textbf{Battle-Tested Insights:} Lessons from years of live trading
\end{itemize}

\textbf{Who Should Read This?}

\begin{enumerate}
    \item \textbf{Aspiring Quants:} Learn systematic trading from the ground up
    \item \textbf{Professional Traders:} Discover cutting-edge techniques and risk management
    \item \textbf{Technology Professionals:} Understand how to build production trading systems
    \item \textbf{Students:} Bridge the gap between academic theory and real-world application
\end{enumerate}

\textbf{The Q23 Philosophy}

Q23 represents the pinnacle of systematic trading technology. Every recipe in this cookbook is drawn from our battle-tested system that manages real capital with:

\begin{itemize}
    \item 24+ production strategies
    \item 40+ sophisticated factors
    \item Real-time risk management
    \item Institutional-grade infrastructure
\end{itemize}

\textbf{Acknowledgments}

This cookbook builds upon the foundational work of Max Dama, who first made quantitative trading accessible. We're honored to extend his legacy with production-grade techniques that honor his spirit while embracing the frontier of 2026 technology.

\textbf{Ready to Start Cooking?}

Let's begin our journey through the most comprehensive quantitative trading cookbook ever created. Each recipe builds upon the last, creating a complete mastery of systematic alpha generation.

\chapter{Foundations: The Quant Kitchen}

\section{Setting Up Your Quant Kitchen}

Before we start cooking, we need to set up our kitchen. The Q23 system provides everything you need for production-grade quantitative trading.

\subsection{Recipe 1.1: Installing Q23 System}

\textbf{Basic Q23 Installation}
\textbf{Ingredients:}
\begin{itemize}
    \item Python 3.9+
    \item Git repository access
    \item Basic data science stack (numpy, pandas, xarray)
\end{itemize}

\textbf{Instructions:}
\begin{enumerate}
    \item Clone the Q23 repository:
\begin{lstlisting}[language=bash]
git clone https://github.com/q23/q23_quant_system.git
cd q23_quant_system
\end{lstlisting}

    \item Install dependencies:
\begin{lstlisting}[language=bash]
pip install -r requirements.txt
\end{lstlisting}

    \item Configure data sources:
\begin{lstlisting}[language=bash]
cp config.template.yaml config.yaml
# Edit config.yaml with your API keys
\end{lstlisting}

    \item Run basic system check:
\begin{lstlisting}[language=bash]
python -c "import q23; print('Q23 system ready!')"
\end{lstlisting}
\end{enumerate}

\textbf{Yield:} Fully functional Q23 quantitative trading system.

\textbf{Pro Tip:} Always run in a virtual environment to avoid dependency conflicts.
}

\subsection{Recipe 1.2: Data Pipeline Setup}

\textbf{Market Data Ingestion Pipeline}
\textbf{Ingredients:}
\begin{itemize}
    \item Market data API credentials (Yahoo, Alpha Vantage, etc.)
    \item Database for storage (PostgreSQL recommended)
    \item Q23 data loader modules
\end{itemize}

\textbf{Instructions:}
\begin{enumerate}
    \item Configure data sources in \texttt{config.yaml}:
\begin{lstlisting}
data_sources:
  yahoo_finance:
    enabled: true
    api_key: "your_key_here"
  alpha_vantage:
    enabled: true
    api_key: "your_key_here"
\end{lstlisting}

    \item Initialize database schema:
\begin{lstlisting}[language=bash]
python -m q23.strategy.data_loader --init-db
\end{lstlisting}

    \item Run initial data download:
\begin{lstlisting}[language=bash]
python -m q23.strategy.data_loader --download --symbols SPY,AAPL,MSFT --start-date 2020-01-01
\end{lstlisting}

    \item Verify data integrity:
\begin{lstlisting}[language=bash]
python -c "from q23.strategy.data_loader import load_market_data; ds = load_market_data(); print(f'Loaded {len(ds.asset)} assets')"
\end{lstlisting}
\end{enumerate}

\textbf{Yield:} Production-ready market data pipeline with 10+ years of historical data.

\textbf{Quality Check:} Ensure no missing values and proper OHLCV structure.
}

\chapter{Mathematical Foundations: The Secret Sauce}

\section{Probability and Statistics: The Base Flavors}

\subsection{Recipe 2.1: Robust Statistical Measures}

\textbf{Q23-Style Robust Statistics}
\textbf{Ingredients:}
\begin{itemize}
    \item Raw financial time series data
    \item Q23 math utilities
    \item Understanding of statistical robustness
\end{itemize}

\textbf{Mathematical Foundation:}
Traditional statistics fail in financial data due to fat tails and outliers. Q23 uses robust statistics:

\begin{align}
\text{Median Absolute Deviation: } &\quad MAD = \mathrm{median}(|X - \mathrm{median}(X)|) \\
\text{Robust Z-Score: } &\quad z_i = \frac{x_i - \mathrm{median}(X)}{MAD \cdot 1.4826 + \epsilon}
\end{align}

\textbf{Implementation:}
\begin{lstlisting}[language=Python,caption=Q23 Robust Statistics]
from q23.shared.math_utils import robust_zscore_xr, safe_corrcoef
import xarray as xr

# Load market data
returns = xr.DataArray(...)  # Your returns data

# Compute robust z-scores across assets
z_scores = robust_zscore_xr(returns, dim="asset", clip=5.0)

# Safe correlation matrix
corr_matrix = safe_corrcoef(returns.values)

print(f"Z-score range: {z_scores.min().values:.2f} to {z_scores.max().values:.2f}")
\end{lstlisting}

\textbf{Yield:} Outlier-resistant statistical measures suitable for financial time series.

\textbf{Key Insight:} Traditional mean/std z-scores can be dominated by single extreme observations. Robust statistics provide more reliable normalization for factor construction.
}

\subsection{Recipe 2.2: Time Series Analysis Toolkit}

\textbf{Advanced Time Series Diagnostics}
\textbf{Ingredients:}
\begin{itemize}
    \item Financial time series (returns, prices, volumes)
    \item Statistical testing framework
    \item Stationarity tests and transformation tools
\end{itemize}

\textbf{Mathematical Foundation:}
Financial time series often exhibit:
\begin{itemize}
    \item \textbf{Non-stationarity:} Unit roots, trends, seasonality
    \item \textbf{Volatility clustering:} GARCH effects
    \item \textbf{Fat tails:} Kurtosis > 3
    \item \textbf{Autocorrelation:} Momentum and mean-reversion
\end{itemize}

\textbf{Implementation:}
\begin{lstlisting}[language=Python,caption=Q23 Time Series Diagnostics]
import numpy as np
import pandas as pd
from statsmodels.tsa.stattools import adfuller, kpss
from q23.shared.math_utils import ewma_vol_1d

def timeseries_diagnostics(series, name="Series"):
    """Comprehensive time series diagnostics."""
    results = {}
    
    # Stationarity tests
    adf_result = adfuller(series.dropna())
    kpss_result = kpss(series.dropna())
    
    results['adf_statistic'] = adf_result[0]
    results['adf_pvalue'] = adf_result[1]
    results['kpss_statistic'] = kpss_result[0]
    results['kpss_pvalue'] = kpss_result[1]
    
    # Distribution properties
    results['mean'] = series.mean()
    results['std'] = series.std()
    results['skewness'] = series.skew()
    results['kurtosis'] = series.kurtosis()
    
    # Autocorrelation
    results['autocorr_1'] = series.autocorr(lag=1)
    results['autocorr_5'] = series.autocorr(lag=5)
    results['autocorr_21'] = series.autocorr(lag=21)
    
    # Volatility clustering
    ewma_vol = ewma_vol_1d(series.values, lam=0.94)
    results['vol_clustering_ratio'] = ewma_vol[-1] / series.std()
    
    return results

# Example usage
spy_returns = load_spy_returns()
diag = timeseries_diagnostics(spy_returns, "SPY Returns")

\end{document}
